\section{Methodology}
The numerical solution of the shallow water equations was implemented using two different grids: an unstaggered grid and a staggered grid. The initial condition for the height field $h$ was defined as a localized cosine wave perturbation given by
\begin{equation}
    h_0(x) = 
    \begin{cases}
        0.5 + 0.5 \cos\left(10 \pi (x-0.5)\right), & 0.4 < x < 0.6, \\
        0, & \text{otherwise.}
    \end{cases}
\end{equation}
And a velocity field $u$ initialized as zero. Both numerical schemes were implemented using the leapfrog time-stepping method, for a total simulation time of $T=1$ arbitrary time units and a spatial domain of $x \in [0,1]$ arbitrary length units, with periodic boundary conditions. Both models used a phase speed of $c=1$, mean fluid depth of $H=1$ and a gravitational acceleration of $g=1$. 

The unstaggered grid scheme was implemented with a spatial resolution of $\Delta x = 0.025$ and a a Courant number of $\mu=0.9$. The staggered grid scheme was implemented with three different spatial resolutions: $\Delta x = 0.05$, $\Delta x = 0.025$ and $\Delta x = 0.0125$, all with a Courant number of $\mu=0.45$. A robert-Asselin filter with a filtering parameter of $\gamma=0.005$ was applied to the staggered grid scheme, while a filter with $\gamma=0.01$ was applied to the unstaggered grid scheme in order to evaluate the effect of the filter on these numerical solutions.

The relative error of the different schemes were computed at each time step by comparing the numerical solution to the analytical solution of the shallow water equations, which is given by 
